\documentclass[aspectratio=169,11pt]{beamer}

\usetheme{Madrid}
\usecolortheme{default}
\usefonttheme{professionalfonts}
\setbeamertemplate{navigation symbols}{}
\setbeamertemplate{footline}[frame number]

\usepackage{amsmath, amssymb, booktabs}

\title{Gendered Policy Incidence}
\subtitle{Childcare, Leave, Taxation, Quotas, Transparency, And Protection\\Gender Study --- Week 6}
\author{Teaching deck draft}
\date{}

\begin{document}

\begin{frame}
  \titlepage
\end{frame}

\begin{frame}{Central question}
\begin{itemize}
    \item Which labor-market margins do gender-related policies move?
    \item Who bears incidence after households, firms, coworkers, and markets adjust?
    \item How do firms respond through pay, staffing, promotion, compliance, classification, and retention?
    \item When does welfare improve, not just the measured gap?
\end{itemize}
\end{frame}

\begin{frame}{Why policy incidence is a labor object}
\begin{itemize}
    \item Policies operate through labor margins: employment, hours, attachment, wages, bonuses, promotion, authority, safety, and sorting.
    \item A policy can reduce one gap while worsening another margin or shifting costs to future careers.
    \item The nominal target is not necessarily the incidence bearer.
    \item Policy variation can reveal deeper gender mechanisms.
\end{itemize}
\end{frame}

\begin{frame}{Bridge from Weeks 2--5}
\small
\begin{tabular}{p{0.25\linewidth} p{0.58\linewidth}}
\toprule
Earlier mechanism & Week 6 stress test \\
\midrule
Household constraints & childcare, leave, tax incentives, care allocation \\
Early choices and sorting & quotas, pipeline policies, exposure to job ladders \\
Firm-side gaps & transparency, equal pay, promotion, bonuses, compliance \\
Norms and bargaining & second-earner incentives, leave take-up, reporting, voice \\
\bottomrule
\end{tabular}
\end{frame}

\begin{frame}{Four objects to keep separate}
\begin{enumerate}
    \item Average treatment effects on outcomes.
    \item Incidence across workers, firms, coworkers, households, and future careers.
    \item Equilibrium and firm response.
    \item Welfare over income, time, safety, autonomy, child outcomes, firm surplus, and fiscal cost.
\end{enumerate}
\end{frame}

\begin{frame}{Policy-to-margin framework}
\[
\tau_m=E\!\left[Y_i^m(P_1)-Y_i^m(P_0)\right]
\]
\begin{itemize}
    \item \(m\): participation, hours, wages, bonuses, promotion, retention, fertility, care, safety, or welfare.
    \item A single policy can have different signs across margins.
    \item Do not interpret a wage-gap effect as a complete welfare effect.
\end{itemize}
\end{frame}

\begin{frame}{Incidence map}
\[
\Delta Y_{ifh}^m
=
\Delta Y^{m,w}_{ifh}
+\Delta Y^{m,h}_{ifh}
+\Delta Y^{m,f}_{ifh}
+\Delta Y^{m,c}_{ifh}
+\Delta Y^{m,eq}_{ifh}
\]
\begin{itemize}
    \item Worker component: preferences, constraints, search, bargaining.
    \item Household component: care allocation, fertility, specialization.
    \item Firm component: pay, staffing, promotion, schedules, compliance.
    \item Coworker and equilibrium components: spillovers, sorting, prices, wages.
\end{itemize}
\end{frame}

\begin{frame}{A welfare object}
\[
W(P)=
\sum_i \omega_i U_i(c_i,h_i,\ell_i,a_i,s_i;P)
+\sum_h \omega_h B_h(P)
+\sum_f \omega_f \Pi_f(P)
+\sum_k \omega_k C_k(P)
-\lambda G(P)
\]
\begin{itemize}
    \item Gap reduction is one input, not the full social objective.
    \item Welfare can rise through safety, autonomy, care quality, or insurance even when wages move little.
    \item Welfare can fall if the measured gap narrows through costly sorting or career penalties.
\end{itemize}
\end{frame}

\begin{frame}{Childcare and care support}
\begin{itemize}
    \item Main margins: participation, hours, job quality, business activity, care time, child outcomes.
    \item Incidence can run through mothers, fathers, children, households, firms, and childcare providers.
    \item Evidence from Quebec and Uganda shows why childcare is a multi-margin intervention.
    \item Labor question: how binding are time constraints and care complementarities?
\end{itemize}
\end{frame}

\begin{frame}{Leave and return to work}
\begin{itemize}
    \item Main margins: attachment, tenure, return timing, experience, promotion, fertility.
    \item Short-run attachment is not long-run career convergence.
    \item Firm responses: replacement hiring, staffing buffers, delayed promotion, project reassignment.
    \item Labor question: are child penalties driven by job continuity, human capital, or firm-side career penalties?
\end{itemize}
\end{frame}

\begin{frame}{Taxation and second earners}
\begin{itemize}
    \item Main margins: participation, hours, formal work, specialization, bargaining.
    \item Joint taxation can tax the secondary earner at a high effective marginal rate.
    \item Responses depend on childcare, norms, household bargaining, and labor demand.
    \item Labor question: what is the elasticity of secondary-earner work, and who captures the gain?
\end{itemize}
\end{frame}

\begin{frame}{Quotas and representation}
\begin{itemize}
    \item Main margins: top representation, authority, networks, pipeline, promotion, spillovers.
    \item Board quotas can change the top node without changing the broader pipeline.
    \item Firm responses: search, mentoring, internal reallocation, compliance appointments.
    \item Labor question: where in the hierarchy is the bottleneck?
\end{itemize}
\end{frame}

\begin{frame}{Transparency and equal pay}
\begin{itemize}
    \item Main margins: base pay, bonuses, raises, negotiation, job classification, retention.
    \item Transparency can narrow gaps through public scrutiny, bargaining, bonus compression, or slower pay growth for men.
    \item Equal-pay rules can induce wage compression or segregation into different job categories.
    \item Labor question: what firm pay-setting rule changed?
\end{itemize}
\end{frame}

\begin{frame}{Protection, reporting, and safety}
\begin{itemize}
    \item Main margins: reporting, retention, mobility, occupational choice, safety, welfare.
    \item Formal protection matters only if workers can credibly use it.
    \item Firm responses: compliance systems, documentation, scheduling, discipline, reporting channels.
    \item Labor question: how do risk, dignity, and reporting costs enter labor supply and job choice?
\end{itemize}
\end{frame}

\begin{frame}{Firm response is central}
\scriptsize
\begin{tabular}{p{0.24\linewidth} p{0.31\linewidth} p{0.31\linewidth}}
\toprule
Policy area & Firm adjustment & Incidence concern \\
\midrule
Leave & staffing, promotion timing, assignment & coworkers and future careers \\
Transparency & pay bands, bonus redesign, raises & lower-paid women vs higher-paid men \\
Equal pay & classification, segregation, compression & target workers vs job reallocation \\
Protection & compliance, reporting, documentation & reporting workers and nonreporters \\
Quotas & search, mentoring, pipeline design & top appointees vs broader cohorts \\
\bottomrule
\end{tabular}
\end{frame}

\begin{frame}{Policy as identification device}
\small
\begin{tabular}{p{0.28\linewidth} p{0.55\linewidth}}
\toprule
Policy shock & Mechanism question \\
\midrule
Childcare expansion & How binding are care-time constraints? \\
Paid leave reform & Is the penalty attachment, experience, or firm response? \\
Tax reform & How elastic is secondary-earner labor supply? \\
Quota & Where does the leadership pipeline bind? \\
Transparency & Are gaps driven by bargaining, bonuses, or pay rules? \\
Equal-pay rule & How do firms adjust when pay discrimination is constrained? \\
Protection reform & How costly is reporting, risk, and unsafe work? \\
\bottomrule
\end{tabular}
\end{frame}

\begin{frame}{Reduced form is not mechanism}
\[
Y_{ift}=\alpha_i+\delta_f+\lambda_t+\beta P_{ft}+X_{ift}'\gamma+\varepsilon_{ift}
\]
\begin{itemize}
    \item \(\beta\) can be a credible policy effect.
    \item It does not reveal whether pay, bonuses, promotions, sorting, or retention moved.
    \item Mechanism claims require additional outcomes or a stronger design.
\end{itemize}
\end{frame}

\begin{frame}{Week 6 lab}
\begin{itemize}
    \item Reproduce: synthetic pay-transparency reporting exercise inspired by Blundell, Duchini, Simion, and Turrell.
    \item Diagnose: classify outputs as outcome effects, incidence, firm response, or welfare-relevant missing margins.
    \item Transfer: equal-pay-for-similar-work exercise inspired by Gentile Passaro, Kojima, and Pakzad-Hurson.
    \item Optional memo: paid leave, board quotas, or childcare subsidies.
\end{itemize}
\end{frame}

\begin{frame}{Reading-heavy frontier week}
\begin{itemize}
    \item Childcare: universal care, legal care regimes, experimental subsidies.
    \item Leave: attachment versus long-run careers.
    \item Taxes: secondary-earner incentives and household allocation.
    \item Quotas: leadership nodes versus pipelines.
    \item Transparency and equal pay: firm pay-setting and regulatory incidence.
    \item Protection: formal rights, enforcement, reporting, and safety.
\end{itemize}
\end{frame}

\begin{frame}{Bridge to Week 7}
\begin{itemize}
    \item Week 6 treats protection and reporting as policy-incidence problems.
    \item Week 7 studies safety, violence, harassment, and mobility as labor-market constraints directly.
    \item Safety changes feasible jobs, reservation wages, sorting, retention, authority, and welfare.
\end{itemize}
\end{frame}

\begin{frame}{Takeaways}
\begin{itemize}
    \item Organize gender policy by margins and incidence, not by policy shopping.
    \item Firm response is part of the treatment.
    \item Average effects, incidence, equilibrium response, and welfare are different objects.
    \item Policy variation can identify mechanisms, but only if the design observes the margin that moves.
\end{itemize}
\end{frame}

\end{document}
